% 原书第 37—39 页；3.4。
\section{结构因果模型的典范表示}
\label{sec:ch03-canonical}

我们已经讨论了结构因果模型能够导出的两类因果陈述：第一类是系统在可能的干预下的行为，第二类是反事实陈述。为了进一步理解二者的区别，下面引入结构因果模型的“\textbf{典范表示（canonical representation）}”。\footnote{文献中多处使用过这种表示，例如 Pearl［2009］；不过，我们没有见到“典范表示”这一名称。}
根据结构赋值
\[
 E=f_E(C,N_E),
\]
对于噪声 $N_E$ 的每一个固定取值 $n_E$，$E$ 都是 $C$ 的确定性函数：
\begin{equation}
 E=f_E(C,n_E).
 \tag{3.8}
\end{equation}
换言之，若 $C$ 和 $E$ 分别在 $\mathcal{C}$ 和 $\mathcal{E}$ 中取值，那么噪声 $N_E$ 就负责在从 $\mathcal{C}$ 到 $\mathcal{E}$ 的不同函数之间切换。因此，不失一般性，可以假设 $N_E$ 的取值空间，就是从 $\mathcal{C}$ 到 $\mathcal{E}$ 的所有函数组成的集合，记为 $\mathcal{E}^{\mathcal{C}}$。采用这一约定，式 (3.8) 也可以改写为
\begin{equation}
 E=n_E(C),
 \tag{3.9}
\end{equation}
我们将它称为联系 $C$ 与 $E$ 的结构方程的典范表示。

接下来说明，为什么两个具有不同典范表示的结构因果模型，可以诱导出相同的干预概率，却给出不同的反事实陈述。为此，我们只考虑 $C$ 在有限集合 $\mathcal{C}=\{1,\ldots,k\}$ 中取值的情形。此时，从 $\mathcal{C}$ 到 $\mathcal{E}$ 的函数集合可以表示为 $k$ 重\textbf{笛卡尔积（Cartesian product）}
\[
 \mathcal{E}^k:=\underbrace{\mathcal{E}\times\cdots\times\mathcal{E}}_{k\text{ 个}},
\]
其中第 $j$ 个分量表示 $C=j$ 时 $E$ 所取的值。相应地，分布 $P_{N_E}$ 就由 $\mathcal{E}^k$ 上的一个联合分布给出；其第 $j$ 个分量的边缘分布，决定了条件分布 $P_{E\mid C=j}$。由于 $C$ 是原因、$E$ 是结果，有 $P_E^{\dox{C:=j}}=P_{E\mid C=j}$。换言之，在这里，干预概率与观测条件概率相同。因此，虽然结构因果模型精确指定了 $P_{N_E}$，也就是所有分量的联合分布，但它在干预方面的因果含义，完全由向量值噪声变量 $N_E$ 的各个分量的边缘分布决定。

与结果变量对应的噪声 $N_E$，其各分量之间的统计相依关系虽然不影响干预因果陈述，却会影响反事实陈述。为说明这一点，令 $C$ 和 $E$ 均为二值变量，即 $\mathcal{C}=\mathcal{E}=\{0,1\}$。从 $\{0,1\}$ 到 $\{0,1\}$ 的函数集合为
\[
 \mathcal{E}^{\mathcal{C}}=\{\mathbf{0},\mathbf{1},\mathrm{ID},\mathrm{NOT}\},
\]
其中，$\mathbf{0}$ 和 $\mathbf{1}$ 分别表示恒取 $0$ 和 $1$ 的常值函数，$\mathrm{ID}$ 表示恒等函数，$\mathrm{NOT}$ 表示取反函数。

为构造两个不同的分布 $P_{N_E}^{1}$ 和 $P_{N_E}^{2}$，使它们诱导出相同的条件分布 $P_{E\mid C=0}$ 和 $P_{E\mid C=1}$，第一种取 $\mathbf{0}$ 与 $\mathbf{1}$ 的等权混合，第二种取 $\mathrm{ID}$ 与 $\mathrm{NOT}$ 的等权混合。两种情况下，$C$ 与 $E$ 都在统计意义上相互独立，对 $C$ 的干预也都不影响 $E$ 的分布，因为无论 $C$ 取何值，$E$ 始终相当于一次公平抛硬币的结果。

在笛卡尔积表示中，这四个函数写作
\[
 \mathcal{E}^{\mathcal{C}}=\{(0,0),(1,1),(0,1),(1,0)\},
\]
其中，第一、第二个分量分别表示 $C=0$ 和 $C=1$ 的像。显然，$(0,0)$ 与 $(1,1)$ 的等权混合，以及 $(0,1)$ 与 $(1,0)$ 的等权混合，在笛卡尔积的第一、第二个分量上，都诱导出相同的边缘分布；这与前面所说的二者诱导出相同干预分布是一致的。然而，反事实陈述“如果将 $C$ 设为另一个值，$E$ 就会取另一个值”，只在 $\mathrm{ID}$ 与 $\mathrm{NOT}$ 的混合下成立，在 $\mathbf{0}$ 与 $\mathbf{1}$ 的混合下则不成立。因此，反事实陈述不仅取决于噪声变量 $N_E$ 各分量的边缘分布，也取决于笛卡尔积各分量之间的统计相依关系。

注意，两个形式不同的结构因果模型，不仅可能诱导出相同的干预分布，甚至还可能导出相同的反事实陈述。对于赋值
\[
 E:=f_E(C,N_E),
\]
显然，对 $N_E$ 进行\textbf{重参数化（reparameterization）}并不影响这些陈述。具体而言，在 $N_E$ 的取值范围上选取某个双射 $g$，并将噪声变量重新定义为 $\widetilde N_E:=g(N_E)$，就可以写成
\[
 E:=\widetilde f_E(C,\widetilde N_E)
   =f_E\bigl(C,g^{-1}(\widetilde N_E)\bigr).
\]
典范表示 (3.9) 消除了这一额外的自由度，避免其干扰我们对反事实的讨论。
